\subsubsection*{The positive cyclotomic valuation}
\label{subsec:positive-twist}

Suppose the cyclotomic minimum \(\nu=\nu_\ell\) of
\eqref{eq:signature-cyclotomic-minimum} is positive. Write \(O=O_0\), choose a
uniformizer \(\pi\), and put
\[ e=\operatorname{ord}_O(\ell^\nu),\qquad I=\ell^\nu\mathfrak m_O=(\pi^{e+1}). \]
The residual calculation detects an entry of the adjugate. One elementary divisor can still carry the entire remaining loss, so the determinant itself must now be detected modulo \(I\).

Let \(b\) be the number of good signatures at this coefficient
place, and \(R_0\) the sum of the bounds for the number of Frobenius elements in their
residual tests, padding each list with identities. Use \(R=R_0+b\) first pairs, and initially
\(n=R_0\) second pairs. Prescribe the first \(R_0\) Frobenius elements of
the second list as in Section~\ref{subsec:residual-twist}.
Future second pairs will use the remaining first coordinates
one at a time. Write \(p=2r\) and set
\[
 w_0=4rR_0,\qquad w_{j+1}=2(e+1)w_j+2+p .
\]
Only the finite signature data enter these constants. Throughout
take \(M\) so large that \(\ell^{M-\nu-1}>w_b+2\) and
\(M\ge\nu+2\).

\begin{lemma}[Independence of low-degree coefficients]
Let \(F,G\) lift the same group-ring element. After reduction to \(O/I\), or to the residue field, their coefficients agree in the range
\[
 [X^\alpha]F=[X^\alpha]G\qquad(|\alpha|\le w_b+1).
\]
\end{lemma}
\begin{proof}
For \(0<i<\ell^M\) one has
\[
 i\binom{\ell^M}{i}=\ell^M\binom{\ell^M-1}{i-1},\qquad
 v_\ell\binom{\ell^M}{i}=M-v_\ell(i),
\]
since \(\binom{\ell^M-1}{i-1}\) is a unit by expansion modulo \(\ell\).
For \(1\le i\le w_b+1<\ell^{M-\nu-1}\) the coefficient belongs
to \(I\). Each coordinate relation \((1+U)^{\ell^M}-1\), and hence
each polynomial multiple, has zero coefficients in these degrees.
\end{proof}

The induction keeps a unit coefficient in a cofactor through
degree \(w_j\) for every signature with a nonempty presentation,
after setting to zero the unused first variables. For each of
the first \(j\) signatures it also keeps a determinant coefficient
outside \(I\) through degree \(w_j\), with all variables allowed.
To start, set the extra first variables to zero. At these places
the residual Frobenius is scalar different from one on either
side, and its norm is one modulo \(\ell\). Writing \(D\) for that residual representation, local duality and the local Euler characteristic give
\[
 H^j(E_v,D)=0\qquad(j=0,1,2).
\]

Successive extensions give the same vanishing on finite
coefficient modules after this specialization. Deleting the
extra pairs therefore gives the old Selmer functor. Its representing
module after scalar change is the old module by naturality.
The first Fitting ideal (the ideal generated by the submaximal
minors of a square presentation) commutes with this base change.
The unit cofactor from Section~\ref{subsec:residual-twist}
therefore supplies the initial assertion for any presentation.
An empty presentation has determinant one; it can also be
stabilized by a unit diagonal entry.

\paragraph{Preserving a coefficient at a larger level.}
Consider a new second pair satisfying the coordinate conditions
and lying in the subgroup \(J\) for every signature.
\begin{lemma}[Fitting ideals at an added pair]
After setting the new second variable to zero, let \(P_+\) and
\(P_0\) be the new and old modules reduced modulo \(I\).
For the Frobenius \(\sigma\) at the new prime put \(u=\eta(\sigma)-1\).
Then
\[
 u^p\operatorname{Fitt}_i(P_0)\subseteq\operatorname{Fitt}_i(P_+)\qquad(i=0,1).
\] A nonzero determinant coefficient modulo \(I\), or
a unit cofactor coefficient modulo \(\pi\) with unused variables
zero, can be retained with degree increased by at most \(p\)
in the asserted low-degree range.
\end{lemma}
\begin{proof}
First set its own variable to zero and reduce modulo \(I\).
Write \(P_+\) for the new representing module after this
specialization and \(P_0\) for the old module. There is a
surjection
\[
 P_+\to P_0
\]
with kernel generated by \(p\)
inertia coordinates. In fact in the new Selmer functor the
unramified classes at the new pair are exactly the old ones.
The old action is unramified there; the only possible
inertia homomorphisms on finite primary modules factor through
the pro-\(\ell\) tame quotient. Evaluating a representative at
one generator on each side gives the claimed rows in \(P_+\),
independent of coboundaries. The kernel description now follows
by testing on all finite modules. Here the specialized rings
are finite.

Put \(u=\eta(\sigma)-1\) for the old universal character on
a Frobenius at the new prime, choosing the direct character
\(\eta\). For the inertia generators \(t_1,\ldots,t_p\) of the kernel this gives the scalar relations
\[
 (\eta(\sigma)-1)t_i=u t_i=0\qquad(1\le i\le p).
\]

Indeed the untwisted Frobenius matrix is one modulo \(I\) on
\(J\) and the tame relation multiplies an inertia value by
the residue norm, which is one in \(O/I\). Inversion on the
conjugate side or in the twist changes this annihilator by
a unit. If \(A_0\) is a presentation matrix for \(P_0\), lift its
relations to the enlarged list of generators for \(P_+\).
Together with the inertia relations these give a block of the
new presentation of the form
\[
        \begin{pmatrix} \widetilde A_0 & *\\ 0 & u I_p\end{pmatrix}
\]
when rows write relations. Adjoining all \(p\) inertia rows
and columns to a minor for \(\operatorname{Fitt}_i(P_0)\)
multiplies it by \(u^p\).

For \(i=0,1\) this proves
\[
        u^p\operatorname{Fitt}_i(P_0)
                \ \subseteq\ \operatorname{Fitt}_i(P_+).
\]
It suffices that these are relations; further relations enlarge
the ideal. One may stabilize either presentation for \(i=1\).
The polynomial for \(u\) has constant term zero and its
linear coefficient at the newly used first variable is one
by the Frobenius coordinate prescription. Even after deleting
the still unused first variables its linear form is monic in
that variable. Write \(U_1\) for its linear part. Over \(O/I\), a monic
polynomial in the newly used first variable is a non-zero-divisor.
If \(F_d\ne0\) is the lowest nonzero homogeneous component of an
element \(F\) in the old ideal, then
\[
 (u^pF)_{d+p}=U_1^p F_d\ne0.
\]

All smaller degrees are zero. The degree test for the coordinate
relations above then shows that any old detection through
degree \(w\) gives some new generator detected through
\(w+p\). Apply this modulo \(I\) to the determinant, and
modulo \(\pi\) to the adjugate with unused first variables
zero. In the latter test the old cofactor survives when the
newly used first variable is restored, since it survived even
with that variable zero.

\end{proof}
\paragraph{Producing a dual class.}
Process one remaining signature, with current bound \(w\).
Put \(D=(e+1)w\). If its old determinant is already detected
through degree \(2D+1\), take a new pair with trivial prescribed
Frobenius modulo the common subgroup and use the preceding
comparison. Otherwise choose an old square presentation with
matrix \(\Phi\) of positive size \(a\) and a polynomial lift
of its determinant with all coefficients through degree
\(2D+1\) in \(I\).

Choose a cofactor, written as an adjugate
entry \(H\), with a unit coefficient through degree \(w\)
after the unused first variables are set to zero. Choose a nonzero homogeneous part \(\bar H_j\) of the residual cofactor in the used variables, with \(j\le w\). The slopes \(\bar a_i\) must satisfy
\[
 \bar a_i\ne0\quad (i\text{ used}),\qquad \bar H_j((\bar a_i)_i)\ne0.
\]
Such a vector exists on the torus over
the algebraic closure and has finitely many algebraic coordinates.
Lift the slopes to units \(a_i\) in a finite unramified DVR
extension \(O'/O\), putting \(a_i=0\) at unused variables.
Work in
\[
          A=(O'/I O')[\tau]/(\tau^{2D+2}),\qquad
          U_i=a_i \tau,\qquad c_0=\ell^\nu \tau^D .
\]
\begin{lemma}[A dual Selmer class in the truncated ring]
For the resulting specialization \(\Psi\) with zero determinant
and unit cofactor coefficient in degree at most \(w\), let
\(H=(\operatorname{adj}\Psi)_{jl}\) denote the selected cofactor.
Column \(l\) gives a primal Selmer class. On the dual side,
\[
 \exists\,y\in S^\perp(A):\qquad c_0y\ne0.
\]
\end{lemma}
\begin{proof}
The coordinate relations hold by the binomial estimate. The
specialized matrix \(\Psi\) has determinant zero. Suppose the
first unit coefficient of the specialized \(H\) occurs at
degree \(d\le w\). Then \(H=\pi B+\tau^d V\) with \(V\) a
unit. Modulo \(H\), the relation \(\tau^d=-\pi B/V\) gives
\(\tau^{(e+1)d}=0\). Thus \(H\mid\tau^D\), since \(D=(e+1)w\).
Because \(\Psi\,\operatorname{adj}\Psi=(\det\Psi)I_a=0\),
column \(l\) represents \(x\in S(A)\).
Writing \(\tau^D=H H'\), one has
\[
 (c_0\tau H)H'=\ell^\nu\tau^{2D+1}\ne0
 \quad\text{in }(O'/IO')[\tau]/(\tau^{2D+2}).
\]

We use the natural identification of kernel vectors with
classes for every coefficient module. Representatives can be
chosen from one cocycle with coefficients in the representing
module, by applying the functor isomorphism to its identity.
A chosen set of lifts of its entries to the free presentation
then gives matrices whose contraction with a kernel vector
represents the corresponding cocycle.

Let \(Y=\operatorname{Hom}(A,\mathbb Q_\ell/\mathbb Z_\ell)\)
be the additive dual with the balanced evaluation pairing.
Write \(S^\perp\) for the Selmer condition on the Tate dual
with inverse character, orthogonal to the signature condition
on \(Y\). The finite-module Selmer order formula gives
\begin{equation}\label{eq:pos-order}
 |\ker\Psi(Y)|=|S^\perp(A)|
             \prod_{v\in Q}|H^0(E_v,T^\vee(1)\otimes\kappa^{-1} A)|.
\end{equation}
Here \(\kappa\) is the specialized inverse universal
character, and the product includes both places of each pair.

To check the local factors for arbitrary finite coefficients,
use an integral basis of the saturated finite-cocycle module
at \(v\mid\ell\). Its residual functions are independent:
a dependency would divide a primitive integral combination
by \(\pi\) in the same lattice. Finitely many evaluations
give an invertible minor. The span on a finite module \(Y\)
thus has order \(|Y|^{r+d_v}\), where \(d_v\) is the de Rham
quotient dimension in Section~\ref{subsec:signature-selmer}.

Let \(Z_v(Y)\) be this span of cocycles, \(B^1_v(Y)\)
its coboundaries, and \(L_v(Y)=Z_v(Y)/B^1_v(Y)\).
The coboundary map has source of order \(|Y|^r\) and
kernel \(H^0_v\), so
\[
 |L_v(Y)|=\frac{|Z_v(Y)|}{|B^1_v(Y)|}
          =\frac{|Y|^{r+d_v}}{|Y|^r/|H^0_v|}
          =|Y|^{d_v}|H^0_v|.
\]
The product of the factors \(|Y|^{d_v}\) cancels the
complex-place factor.
At the full local conditions away from \(\ell\) the local
Euler characteristic leaves the dual invariants. They vanish
at the exceptional places by the inertia test. Outside the
ramification set unramified conditions are paired. Both global
invariant groups vanish by the scalar marking. This proves
\eqref{eq:pos-order}, and the same proof applies to the
perfectly paired modules \(Y/c_0\tau Y\) and \(A[c_0\tau]\).
\cite{milne-duality}

There exists \(y\in S^\perp(A)\) with \(c_0[y]\ne0\).
Indeed at a first pair the dual invariants vanish by the
residual Frobenius test and finite dévissage. At each current
second pair the inertia of its own coordinate has value
\(1+a_i \tau\) with \(a_i\) a unit. For any invariant \(v\) at such a place,
\[
 (1+a_i\tau)v=v\quad\Longrightarrow\quad \tau v=0 .
\] Suppose all the dual classes were killed
by \(c_0\). For a representative \(y\), write
\[
 c_0y=dm.
\]

The image of \(m\) in the coefficient quotient by \(c_0\)
is invariant globally, hence zero. Subtracting a coboundary
from \(y\) gives values in \(A[c_0]\). At a strict place a
local expression \(y=dm_v\) then has \(c_0m_v\) invariant,
so \(c_0\tau m_v=0\). At \(\ell\) every finite local class
on \(Y/c_0\tau Y\) lifts from \(Y\) by the local description;
naturality of the pairing gives the required orthogonality.
The class thus lifts to the dual Selmer group on \(A[c_0\tau]\).
The invariant products in \eqref{eq:pos-order} coincide,
so
\[
 |\ker\Psi(Y)|\le |\ker\Psi(Y/c_0\tau Y)|.
\]

For an endomorphism \(U\) of a finite module \(M\),
\[
 |\ker U|=|M|/|U(M)|=|\operatorname{coker}U|.
\] Accordingly, the kernel inequality implies
that the quotient surjection
\[
 \frac{Y^a}{\Psi Y^a}
 \longrightarrow\frac{Y^a}{\Psi Y^a+c_0\tau Y^a}
\]
is bijective. Hence \(c_0\tau Y^a\subseteq\Psi Y^a\). Applying
row \(j\) of \(\operatorname{adj}\Psi\), which annihilates this
image, to \(c_0\tau Y\) in coordinate \(l\) shows
\((c_0\tau H)Y=0\). The evaluation pairing of \(A\) with \(Y\)
is perfect, implying \(c_0\tau H=0\), a contradiction.

\end{proof}
\paragraph{A Frobenius element with the prescribed evaluations.}
Intersect \(J\) with the kernels of the used coordinates to
obtain \(J'\). This subgroup is normal and conjugation-stable.
The two twisted actions are trivial on it, the commutators
\([g,z]\) lie in it, and \(\kappa(z)=1\) by the ramification
condition. Both restrictions of the cocycles to \(J'\) are
therefore additive, independent of representatives and
equivariant under conjugation.

\begin{lemma}[A second-half Frobenius with nonzero pairing]
For the primal and dual classes just chosen, there is a new prime
pair, avoiding the previous exceptional set, with the prescribed first-coordinate values, satisfying at the chosen prime \(P\), with Frobenius \(\sigma\),
\[
 \mathrm NP\equiv1\pmod{2\ell^M},\qquad B_\sigma\ne0.
\] Its untwisted Frobenius and that
of its conjugate act trivially modulo \(I\).
\end{lemma}
\begin{proof}
Some \(h\in J'\) has
\[
 B_h=x(h)\cdot y(h)-x(h^c)\cdot y(h^c)\ne0\quad\text{in } A .
\]
For suppose all these differences vanished. Polarizing at
\(h,k\in J'\) gives equality of the two cross sums. Replacing
\(h\) by \(zhz^{-1}\) multiplies its \(x\)-value by \(d\)
and its \(y\)-value by \(d^{-1}\), with the inverse multipliers
on the conjugate pair. Since \(d-d^{-1}\) is a unit, subtraction
separates the two terms and yields
\[
            x(h)\cdot y(k)=x(k^c)\cdot y(h^c).
\]

Let \(R\) and \(R'\) be the twisted primal and dual actions.
Conjugate duality and anti-invariance of \(\kappa\) identify
\((R')^c\) with \(\chi R\), \(\chi=\epsilon^{\text{tw}}\) with \(\text{tw}\) the
signature's fixed twist exponent.
Replace \(h\) by \(ghg^{-1}\) in the last identity.
Equivariance on \(J'\) gives, for every \(g\),
\[
 \operatorname{tr}(R(g)U)=\operatorname{tr}(\chi(g)R(g)V),
 \qquad U=x(h)y(k)^{\mathsf T},
\]
where \(V\) is the conjugate of \(y(h^c)x(k^c)^{\mathsf T}\)
by the fixed intertwiner. The span of \(R(g)\) is the full
matrix algebra over \(A\), by residual absolute irreducibility
and Nakayama. The same holds for \(\chi R\).

Consider the subalgebra spanned by the pairs \((R(g),\chi(g)R(g))\) in
\(M_r(A)^2\). It surjects onto both factors. Its intersections
with the separate factors are two-sided ideals, of the form
\(M_r(I_1)\) and \(M_r(I_2)\) by multiplication with matrix
units. If \(\mathcal A\) is this subalgebra, a pair \((B,C)\in\mathcal A\)
associates the class of \(B\) modulo \(M_r(I_1)\) to the class
of \(C\) modulo \(M_r(I_2)\). Surjectivity and the definitions
of \(I_1,I_2\) give a well-defined inverse in the other direction.
It is an isomorphism
\[
      M_r(A/I_1)\xrightarrow{\sim} M_r(A/I_2)
\]
of \(A\)-algebras, so their scalar annihilators coincide and
\(I_1=I_2=I'\).

If the quotient is nonzero its matrix algebra
automorphism is inner. Indeed the images of the diagonal matrix
units cut the standard module over the local ring \(A/I'\)
into isomorphic projective summands, necessarily free of rank
one. Transport a generator of the first by the images of the
off-diagonal matrix units; the resulting basis conjugates
all matrix units at once. Taking determinants in the quotient
gives \(\chi(g)^r-1\in I'\) for every \(g\), also true for the
zero quotient. By definition of \(\nu\), these scalars generate
\((\ell^\nu)\) after coefficient extension. The trace identity
on \((\ell^\nu E_{ab},0)\) gives
\[
 \ell^\nu x(h)_i y(k)_j=0
\]
for all coordinates.

Let \(J_x\) be the ideal generated by the entries of all
\(x(h)\), \(h\in J'\). Modulo \(J_x\) the cocycle factors
through \(G_E/J'\), on which \(z\) is central and acts by
a scalar differing from one by a unit. Comparing the cocycle law at \(zg\) and \(gz\) gives
\[
        (d-1)x(g)=(g-1)x(z)
\]
over \(A/J_x\). As \(d-1\) is a unit, this expresses the
cocycle as a coboundary with constant \((d-1)^{-1}x(z)\). Naturality and injectivity
of the kernel description put every entry of column \(l\)
of the adjugate in \(J_x\).

In particular \(\tau^D\in J_x\)
since \(H\mid\tau^D\). As \(\ell^\nu J_x y(k)=0\) for \(k\in J'\),
this gives
\[
 c_0y|_{J'}=0.
\]
In the quotient
\(G_E/J'\), the central element \(z\) acts on the dual by \(d^{-1}\)
with \(d^{-1}-1\) invertible. For \(f=c_0y\) the two evaluations of \(f(zg)=f(gz)\) give
\[
 (d^{-1}-1)f(g)=(g-1)f(z).
\]
Dividing by the unit exhibits \(f\) as a coboundary, contradicting the preceding lemma.

Choose an open normal subgroup in \(J'\) on which both
cocycles and their conjugates vanish. This is possible since
their restrictions are continuous homomorphisms to finite
modules; intersect the kernels and take the normal core.
This subgroup can contain all inertia at the unused first
pairs. There the used characters and signature actions are
unramified and the residual Frobenius test gives vanishing
of local \(H^1\) for both modules by finite dévissage.
Local coboundaries vanish on inertia, also at every conjugate
place and for every conjugate of its inertia group.
By Chebotarev choose a new second pair with Frobenius
\(\sigma\) conjugate to \(h\) modulo this subgroup, prescribing the first-coordinate vector
\[
 (0,\ldots,0,1,0,\ldots,0)
\]
with the single \(1\) in the next unused first coordinate (additive notation). Indeed \(h\) is in
the kernels of the used coordinates and the unused first
inertia groups can correct all the remaining coordinates.

All old used characters then vanish on the new decomposition
group. The Frobenius can be adjusted by its own inertia to
give value zero also in the new coordinate. Equivariance and
the dual pairing multiply \(B_h\) under conjugation by a
cyclotomic unit, equal on the two conjugate terms. Thus
\(B_\sigma\ne0\). The new prime avoids all exceptional
places, has norm one modulo \(2\ell^M\), and its untwisted
action and that at its conjugate are trivial modulo \(I\).
At each stage the open subgroup and avoidance set may include
the common earlier Frobenius requirements of the other signatures.

\end{proof}
\begin{proposition}[A determinant coefficient after all positive-valuation steps]\label{prop:positive-signature-determinant}
Starting from Lemma~\ref{lem:residual-selmer-coefficient},
under the hypotheses in this construction, after choosing the
reserve pairs, there is a common degree bound \(D\) such that every signature determinant \(F_s\) has a coefficient satisfying
\[
 \exists\alpha:\quad |\alpha|\le D,\qquad v_\ell([X^\alpha]F_s)\le\nu_\ell.
\] The number of pairs, degree and depth threshold depend
only on the original finite data. For each sufficiently large
\(M\), one choice of pairs serves all signatures. If
\(\nu_\ell=0\), no reserve pair is needed.
\end{proposition}

\begin{proof}
\proofstep{Make the new coordinate substitutions before taking the Fitting ideal.}
Use the finite ring
\[
 B=A[X,Y]/(X^2,Y^p),
\]
over \(A\) by
constant inclusion and with projection \(B\to A\).
Keep the old used variables \(a_i \tau\), set the reserve first
variable corresponding to this pair to \(Y\), and set its
own second variable to \(X\). These assignments satisfy all
coordinate relations by the binomial bound. At the new
prime take a tame generator \(\gamma\) mapped to one in the
new coordinate. The inverse universal character has values
\((1+X)^{-1}\) and \((1+Y)^{-1}\) on \(\gamma,\sigma\);
the conjugate pair has inverse values. On projection to \(A\)
all these values become one.

Let \(P'\) be the new representing module after this base
change. For its universal cocycle write \(u_i,t_i\in P'\)
for the Frobenius and inertia entries at the two places,
\(1\le i\le p\). The tame relation gives
\[
          X u_i=Y\eta_i t_i,\qquad
             \eta_i\in B^\times,\quad \eta_i\bmod(X,Y)=1 .
\]
To verify it write \(q\) for the norm. On the pro-\(\ell\)
tame quotient
\[
 \sigma\gamma=\gamma^q\sigma.
\]

For an inertia scalar \(1+v\), \(v^2=0\), the cocycle on
\(\gamma^q\) is \((q+\binom q2 v)\) times its value on
\(\gamma\), hence just its value on \(\gamma\), since
\(q\equiv1\bmod 2\ell^M\). Wild and prime-to-\(\ell\) tame
inertia in the kernel of the action have zero cocycle values.
Expanding the relation now gives the displayed identity,
with
\[
 \eta_i=(1+Y)^{-1}
\]
or one.

In \(P'/(X,Y)P'\), quotienting by the inertia entries gives
the old module with presentation \(\Psi\), by the unramified
condition. Choose lifts of its \(a\) generators. The projected
Frobenius rows modulo inertia are the old rows \(U_i\)
evaluated at \(\sigma,\sigma^c\); an ambiguity of the
representative is a coboundary and vanishes at these elements
of trivial action. Define \(k_i\) to be the entries of
\(y(\sigma)\) on the first side and \(-y(\sigma^c)\) on the
second. Global reciprocity gives a further relation
\[
                \sum k_i t_i=0 \pmod{(X,Y)P'} .
\]
Indeed test against the additive dual of \(A\). A new relaxed
class on that dual with the old character pairs with the old
dual class \(y\). The local cup invariants vanish at the old
places by their conditions, and outside the level by
unramifiedness.

At each of the two new places the action is
trivial and \(y\) is unramified. Only the new inertia value pairs with \(y\)'s Frobenius value. To determine the two signs, identify the scalar Tate coefficients by a primitive \(\ell^{\nu+1}\)-st root. By local reciprocity the unramified character taking Frobenius to one pairs to a generator with the Kummer class of a uniformizer, whose tame pro-\(\ell\) inertia value is a unit. Unramified parts pair to zero.
Conjugation takes the chosen primitive root to its inverse. Thus, in the additive identification of the scalar Tate coefficients,
\[
 \zeta^a\longmapsto\zeta^{-a},\qquad a\longmapsto-a.
\]
This gives the opposite sign at the conjugate place. Applying the
sum of invariants and varying the homomorphism to the additive
dual proves the relation, since that dual is a cogenerator
for finite \(A\)-modules. \cite{milne-duality,milne-cft}

The lifted old generators and \(\eta_i t_i\) generate \(P'\)
by nilpotence of \((X,Y)\). Use them in this order. In the
reduction to \(A\), write the Frobenius rows as \((U_i,*)\),
and lift to rows over \(B\). The top \(a\) relations may be
chosen to reduce to \((\Psi,C)\) for some \(C\), except that
row \(l\) is replaced by \((0,k)\) by reciprocity. Each
reduced relation lifts: its sum lies in \((X,Y)P'\), so
subtract an expression in the same generators with
coefficients in \((X,Y)\).

The additional \(p\) exact tame rows are
\[
 Y(0,E_i)-X(\widetilde U_i,*)\qquad(1\le i\le p),
\]
where \(E_i\) is the \(i\)-th coordinate row. The determinant of this square
array lies in \(\operatorname{Fitt}_0(P')\). Expand in the
last \(p\) rows. The terms with no \(X\) or with at least two copies vanish, since \(Y^p=X^2=0\). For the surviving terms
\[
 X\,(XY^{p-1})=0,\qquad Y\,(XY^{p-1})=0,
\]
so all their other coefficients
may be reduced to \(A\). Expanding now on row \(l\) and the
\(p-1\) unchanged coordinate rows gives, up to one common sign,
\[
       XY^{p-1}\sum_i k_i(U_i\operatorname{adj}\Psi)_l
                      =XY^{p-1}B_\sigma\ne0.
\]
The last inequality uses the monomial basis over \(A\).

\proofstep{Use the bounded nilpotence degree to locate a coefficient.}
Consequently the specialized new determinant, which generates
this Fitting ideal, is nonzero. A monomial of total degree
greater than \(2D+1+p\) vanishes on this substitution:
the old variables are multiples of \(\tau\), and the two others
have exponents at most one and \(p-1\). Thus every lift of
the determinant has a coefficient outside \(I\) in the
asserted degree range. This reasoning applies to any new
square presentation for the same functor.

\proofstep{Use the common degree bound while preserving earlier signature detections.}
The final degree bound must retain the coefficients already detected at earlier signatures.

Repeat this construction for the \(b\) signatures in turn.
The Fitting comparison at each new pair preserves all earlier
detections and the available cofactors within the next bound
\(w_{j+1}\). For each signature determinant \(F_s\), the final common degree bound reads
\[
 \min_{|\alpha|\le w_b} v_\ell([X^\alpha]F_s)\le\nu_\ell.
\]
Here the multi-index ranges over all the variables. For each sufficiently large
\(M\), the same choice of pairs serves all signatures, and the
pairs may depend on \(M\). For \(\nu_\ell=0\)
the residual argument finishes without the reserve pairs.
The number of pairs, the degree and the depth threshold may
depend on the original finite data at the coefficient place;
they are fixed before increasing \(M\). The subsequent
specialization of Section~\ref{subsec:retained-twist} retains
the valuation allowance \(\nu_\ell+\delta\).
\end{proof}

